\section{Week 7: Mode-2 Signature Codec}

The main Week~7 lane moved to mode-2 signature serialization. The actual
standalone wrappers \texttt{pack\_sig\_mode2\_full\_jazz} and
\texttt{unpack\_sig\_mode2\_full\_jazz} were regenerated and pinned by hash.
Their reachable closures contain the actual prefix, full, and rANS helper
procedures used below.

\subsection*{Audited parameters and 1474-byte layout}

The extracted wrapper calls agree on:
\[
\begin{array}{l}
\mathsf{sigbytes}=1474,\;
\mathsf{lcount}=4,\;
\mathsf{hb\_count}=1024,\;
\mathsf{hb\_m}=13,\;
\mathsf{hb\_offset}=6,\\
\mathsf{h\_count}=512,\;
\mathsf{h\_m}=13,\;
\mathsf{h\_offset}=239,\\
\mathsf{base\_hb}=132,\;
\mathsf{base\_h}=7,\;
\mathsf{payload\_limit}=416.
\end{array}
\]

The byte structure recovered from the actual extracted code is:
\begin{center}
\begin{tabular}{r l}
$[0,32)$ & 256 challenge bits, packed within 32 bytes \\
$[32,1056)$ & 1024 signed low-coefficient bytes \\
$[1056,1058)$ & hb/h encoded-size deltas \\
$[1058,1058+\mathsf{hbsize})$ & hbz encoded payload \\
following bytes & h encoded payload \\
through byte 1473 & canonical zero padding
\end{tabular}
\end{center}
The arithmetic boundary is $1056+2+416=1474$. The unpacker additionally
checks nonzero size values, the two base-offset ranges, their total against
416, and zero trailing bytes.

\subsection*{Actual prefix procedures}

The authored theory first proves two procedure-level layout theorems. The
packer theorem opens \texttt{\_pack\_sig\_prefix} and establishes the
challenge-bit and 1024-byte low layouts. The unpacker theorem opens
\texttt{\_unpack\_sig\_prefix} and establishes the inverse bit extraction and
signed-byte extension. Its postcondition also preserves low-array words
$[1024,2048)$.

The sequential harness calls the two actual generated procedures:
\begin{verbatim}
sig <@ Pack._pack_sig_prefix(sig, cp, low, W64.of_int 4,
                             W64.of_int 1474);
(decoded_cp, decoded_low) <@
  Unpack._unpack_sig_prefix(decoded_cp, decoded_low, sig,
                            W64.of_int 4);
\end{verbatim}

The fresh-compiled theorem is
\path{pack_unpack_sig_prefix_mode2_roundtrip}. Its exact semantic
statement is:
\[
\begin{split}
 &\mathsf{canonical\_challenge}(cp)\;\land\;
   \mathsf{canonical\_signed\_low}(low)\\
 &\quad\Longrightarrow
   \mathsf{challenge\_prefix\_eq}(decoded\_cp,cp)\;\land\\
 &\qquad\qquad
   \mathsf{low\_mode2\_eq}(decoded\_low,low)\;\land\;
   \mathsf{low\_tail\_frame}(initial\_low,decoded\_low).
\end{split}
\]

Here challenge words are exactly zero or one. A canonical low word has signed
value in $[-128,128)$ and equals the arithmetic sign extension of its low
byte. The theorem is Hoare partial correctness; it does not assert signing
loop termination. The compiled
\path{prefix_codec_preconditions_satisfiable} lemma supplies all-zero
arrays, establishing non-vacuity.

\subsection*{Proof strategy and failed decomposition}

The pack proof uses three invariants: zero-fill progress, bit-by-bit assembly
of each challenge byte, and the low-byte store cursor. The unpack proof uses a
nested byte/bit invariant for challenge recovery and a signed-extension
invariant for low coefficients. The initial composition attempt fixed the
signature to an external logical parameter. It failed because the local
program variable returned by the preceding pack call could not instantiate
that pure parameter in a \texttt{call} tactic. Generalizing the unpack theorem
to predicates over its current \texttt{sigp} argument closed the composition;
repeating the fixed-array or abstract-function approach should be avoided.

An attempted monolithic pack tail-frame proof was not retained. The compiled
boundary proves the unpack low-array tail, while the pack signature-array
$[1474,2948)$ frame and challenge-output tail remain
\texttt{OBL-SIG-PREFIX-FRAME}. The standalone codec procedures do not write
global memory.

\subsection*{Suffix obligations and security boundary}

The actual metadata and zero-padding gates are audited but do not yet have
procedure theorems. Both rANS pairs remain open:
\begin{itemize}
\item \path{_encode_hb_z1_full} with \path{_decode_hb_z1_full};
\item \path{_encode_h_full} with \path{_decode_h_full}.
\end{itemize}

The prefix result is security-relevant. Deterministic canonical parsing
removes one source of signature ambiguity and malleability. It supplies one
input to the adapter between implementation acceptance and the paper forgery
event. It does not make the full encoding loss zero: suffix ambiguity,
malformed rejection, norm, hint, arithmetic, and challenge consistency remain
open. Consequently neither Sign-output tail reach nor full Verify correctness
follows from this chapter.
